\section*{अध्याय \ref{ch:b3:lithosphere-thermal-models} --- स्थलमंडल के ऊष्मीय मॉडल}

\begin{solution}{exo:b3:lithosphere-thermal-models:1}
$L^2/\kappa = \qty{e12}{s}$, \qty{e14}{s} और \qty{e16}{s}: 0.032, 3.2 और
\qty{317}{Myr}। एक किलोमीटर अपना तापमान दसियों हज़ार
वर्षों में भूल जाता है, सौ किलोमीटर सैकड़ों मिलियन वर्षों में।
\end{solution}

\begin{solution}{exo:b3:lithosphere-thermal-models:2}
$u = \qty{e-3}{m}/\qty{3.156e7}{s} = \qty{3.17e-11}{m.s^{-1}}$, अतः
$\mathrm{Pe} = 3.17\times10^{-11}\times\num{3e4}/10^{-6} = 0.95$। अभिवहन और
चालन समान महत्व के हैं: ऊपर उठती चट्टान ऊष्मा को ऊपर ले जाती है
और सतह के निकट जियोथर्म को तीव्र बनाती है, जबकि केवल चालन प्रवणता को
कम करके आँकेगा।
\end{solution}

\begin{solution}{exo:b3:lithosphere-thermal-models:3}
$\sqrt{\kappa t} = \sqrt{\num{0.8e-6}\times5\times\num{3.156e13}} =
\qty{11236}{m}$, $\eta = 10000/(2\times11236) = 0.445$, और $T = 1450\times0.471
= \qty{683}{\celsius}$।
\end{solution}

\begin{solution}{exo:b3:lithosphere-thermal-models:4}
$490/\sqrt1 = 490$, $490/5 = 98$ और $490/10 = \qty{49}{mW.m^{-2}}$। संवेदक
सबसे संभावित रूप से पहले को कम आँकेगा: नवीन भूपर्पटी दरारयुक्त और
परिसंचारी समुद्री जल के लिए खुली है, जो गर्म सोतों द्वारा ऊष्मा का कुछ भाग दूर ले जाता है।
\end{solution}

\begin{solution}{exo:b3:lithosphere-thermal-models:5}
$z = 2\times1.163\sqrt{\kappa t}$: जहाँ $\sqrt{\kappa t} = 15.89$, 31.78 और
\qty{50.25}{km}, $z = 37.0$, 73.9 और \qty{116.9}{km}। शीतल परत
\qty{10}{Myr} से \qty{40}{Myr} तक दोगुनी होती है; \qty{100}{Myr} तक यह
\qty{95}{km} मोटी प्लेट से पहले ही अधिक गहरी हो जाती है।
\end{solution}

\begin{solution}{exo:b3:lithosphere-thermal-models:6}
$C$, $\alpha$ के समानुपाती है: $\alpha = \num{3.1e-5}\times350/370.5 =
\qty{2.93e-5}{K^{-1}}$। $C$ भी $T_m - T_0$ के समानुपाती है (प्रावार का तापमान
$1450\times350/370.5 = \qty{1370}{\celsius}$) और $\sqrt\kappa$ के
समानुपाती है (विसरणता $0.8\times(350/370.5)^2 = \qty{0.71e-6}{m^2.s^{-1}}$): केवल गहराई
इन दोनों को अलग नहीं कर सकती, इसी कारण ऊष्मा-प्रवाह भी प्रयोग में लाया जाता है।
\end{solution}

\begin{solution}{exo:b3:lithosphere-thermal-models:7}
$3.1\times1450/\num{95e3} = \qty{47.3}{mW.m^{-2}}$ और $3.1\times1350/\num{125e3}
= \qty{33.5}{mW.m^{-2}}$। अर्ध-समष्टि इन मानों तक $(490/47.3)^2 =
\qty{107}{Myr}$ और $(490/33.5)^2 = \qty{214}{Myr}$ पर पहुँचता है: पतली, गर्म प्लेट
पुराने तल के ऊष्मा-प्रवाह को लगभग \qty{100}{Myr} से \qty{47}{mW.m^{-2}} पर बनाए रखती है,
जबकि मोटी, ठंडी प्लेट इसे
\qty{40}{mW.m^{-2}} से नीचे गिरने देती, जो \qty{150}{Myr} पर अर्ध-समष्टि का मान है।
\end{solution}

\begin{solution}{exo:b3:lithosphere-thermal-models:8}
$A_0h_r = 0.4\times60 = \qty{24}{mW.m^{-2}}$, अतः $A_0 =
\num{24e-3}/\num{e4} = \qty{2.4}{\micro W.m^{-3}}$ और $q_m = 60 - 24 =
\qty{36}{mW.m^{-2}}$। \qty{35}{km} पर, $T = \num{36e-3}\times\num{35e3}/2.5 +
\num{2.4e-6}\times\num{e8}/2.5\times(1 - e^{-3.5}) = 504 + 96\times0.970 =
\qty{597}{\celsius}$।
\end{solution}

\begin{solution}{exo:b3:lithosphere-thermal-models:9}
$h_r = \Delta q_0/\Delta A_0 = \num{18e-3}/\num{1.8e-6} = \qty{10}{km}$, और
$q_m = q_0 - h_rA_0 = 66 - 30 = \qty{36}{mW.m^{-2}}$, क्योंकि $h_rA_0 =
\num{e4}\times\qty{3.0e-6}{W.m^{-3}} = \qty{30}{mW.m^{-2}}$।
\end{solution}

\begin{solution}{exo:b3:lithosphere-thermal-models:10}
(क) $Q(t) = \int_0^tk(T_m - T_0)/\sqrt{\pi\kappa t'}\,\mathrm dt' =
2k(T_m - T_0)\sqrt{t/(\pi\kappa)}$। (ख) खोई गई ऊष्मा स्तंभ से हटाई गई
ऊष्मा-अंतर्वस्तु है, $Q = \rho_mc\int_0^\infty(T_m - T)\,\mathrm dz$, और
समस्थिति देती है $(\rho_m - \rho_w)(d - d_r) = \rho_m\alpha\int_0^\infty(T_m -
T)\,\mathrm dz = \alpha Q/c$। (ग) $c = 3.1/(3300\times\num{0.8e-6}) =
\qty{1174}{J.kg^{-1}.K^{-1}}$; \qty{40}{Myr} पर, $t = \qty{1.262e15}{s}$ और
$Q = 2\times3.1\times1450\times\sqrt{\num{1.262e15}/(\pi\times\num{0.8e-6})} =
\qty{2.015e14}{J.m^{-2}}$; $d - d_r = \num{3.1e-5}\times\num{2.015e14}/
(1174\times2270) = \qty{2343}{m}$, जो $370.5\sqrt{40}$ के बराबर है।
\end{solution}

\begin{solution}{exo:b3:lithosphere-thermal-models:11}
$n = 3$ का पद $3104\times(8/\pi^2)\times\tfrac19e^{-9t/36.2} =
280\,e^{-9t/36.2}$ मीटर है: \qty{218}{m} (\qty{1}{Myr} पर), \qty{23}{m} (
\qty{10}{Myr} पर) और \qty{1.9}{m} (\qty{20}{Myr} पर)। यह \qty{10}{m} तक गिरता है जब
$e^{-9t/36.2} = 10/280$ हो, $t = (36.2/9)\ln28 = \qty{13.4}{Myr}$ पर।
\end{solution}

\begin{solution}{exo:b3:lithosphere-thermal-models:12}
(क) $z \gg h_r$ के लिए, $e^{-z/h_r} \to 0$। (ख) $q_m = \qty{24}{mW.m^{-2}}$,
$A_0h_r^2/k = \num{16e-3}\times\num{e4}/2.5 = \qty{64}{K}$, और $z = (1300 -
64)\times2.5/\num{24e-3} = \qty{128.8}{km}$। (ग) $q_m = 2.5\times1236/\num{250e3}
= \qty{12.4}{mW.m^{-2}}$: \omterm{def:b3:lithosphere-thermal-models:reduced}{न्यूनीकृत ऊष्मा-प्रवाह} लगभग आधा होना चाहिए, जैसा
किसी क्रेटन के नीचे होता है, जिसकी मोटी, ठंडी मूल नीचे से कम ऊष्मा आने देती है।
\end{solution}

\begin{solution}{pb:b3:lithosphere-thermal-models:1}
\textbf{1.} $\sqrt{\num{0.8e-6}\times\num{3.156e14}} = \qty{15889}{m}$ और
$\sqrt{\num{0.8e-6}\times\num{1.2624e15}} = \qty{31778}{m}$।

\textbf{2.} $\eta = 20000/31778 = 0.6294$ और $20000/63556 = 0.3147$:
$T = 1450\times0.6266 = \qty{909}{\celsius}$ (\qty{10}{Myr} पर) और
$1450\times0.3437 = \qty{498}{\celsius}$ (\qty{40}{Myr} पर)।

\textbf{3.} $z = 2\times0.7174\sqrt{\kappa t}$: \qty{22.8}{km} और
\qty{45.6}{km}। यह \qty{60}{km} तक पहुँचता है जब $\sqrt{\kappa t} = 60/1.4348 =
\qty{41.82}{km}$, $t = (\num{41.82e3})^2/\num{0.8e-6} = \qty{2.186e15}{s}$ पर,
\qty{69.3}{Myr}।

\textbf{4.} $q = 3.1\times1450/(\sqrt\pi\times15889) = \qty{159.6}{mW.m^{-2}}$ (\qty{10}{Myr} पर) और
\qty{79.8}{mW.m^{-2}} (\qty{40}{Myr} पर); $C_q = 159.6\sqrt{10}
= \qty{505}{mW.m^{-2}.Myr^{1/2}}$।

\textbf{5.} $u = \qty{7.92e-10}{m.s^{-1}}$, $\mathrm{Pe} = 7.92\times10^{-10}
\times\num{95e3}/\num{0.8e-6} = 94$। ऊष्मा प्लेट के साथ पार्श्व दिशा में
उसके पार्श्व विसरण की तुलना में सौ गुना तेज़ चलती है, इसलिए प्रत्येक
स्तंभ केवल ऊर्ध्वाधर चालन द्वारा ठंडा होता है।

\textbf{6.} $(\num{95e3})^2/\num{0.8e-6} = \qty{1.13e16}{s}$, \qty{357}{Myr}।
\qty{150}{Myr} पर, $\sqrt{\kappa t} = \qty{61.5}{km}$ और ठंडी परत,
उसकी कुछ गुना मोटी परत \qty{95}{km} से काफ़ी नीचे तक पहुँचती है:
इतनी पुरानी महासागरीय तली के लिए अर्ध-समष्टि मॉडल मान लेता है कि ठंडापन उन गहराइयों तक पहुँच चुका है जहाँ प्रावार संवहन करता है।

\textbf{7.} $(\rho_m - \rho_w)(d - d_r) = \rho_m\alpha\int(T_m - T)\,\mathrm dz$,
कटक स्तंभ और पुराने स्तंभ के द्रव्यमान प्रतिपूर्ति गहराई के ऊपर
बराबर होते हैं। $T_m - T = T_m\operatorname{erfc}\eta$ का मान
$T_m$ के \qty{0.5}{\percent} से नीचे गिर जाता है, $\eta > 2$ के लिए, इसलिए गहरा भाग कुछ नहीं जोड़ता।

\textbf{8.} $C = 2\times3300\times\num{3.1e-5}\times1450/2270\times
\sqrt{\num{0.8e-6}\times\num{3.156e13}/\pi} = 0.13069\times2834.8 =
\qty{370.5}{m.Myr^{-1/2}}$; $d = 2500 + 370.5\sqrt{10} = \qty{3672}{m}$ और
$2500 + 370.5\sqrt{40} = \qty{4843}{m}$।

\textbf{9.} मॉडल $2500 + 370.5\sqrt{42.5} = \qty{4915}{m}$ देता है,
\qty{440}{m} माध्यिका से गहरा है। माध्यिका में अवसाद का भार,
समुद्री पर्वत और ज्वालामुखीय पठार, हॉटस्पॉट के उभार, और
विभंग क्षेत्रों तथा किनारों के पार्श्वों के स्तंभ शामिल हैं, जो सभी तली को
उथला बना देते हैं।

\textbf{10.} $C$ का मान $T_m$ के समानुपाती है: $370.5\times1350/1450 =
\qty{344.9}{m.Myr^{-1/2}}$, जो समंजित 350 के निकट है।

\textbf{11.} दरारों वाली युवा भूपर्पटी में परिसंचारी समुद्री जल ऊष्मा का
कुछ भाग, संवेदकों के बीच, गर्म सोतों के माध्यम से बाहर ले जाता है।

\textbf{12.} $\tau = (\num{95e3})^2/(\pi^2\times\num{0.8e-6}) =
\qty{1.143e15}{s} = \qty{36.2}{Myr}$; आयाम है $3300\times
\num{3.1e-5}\times1450\times\num{95e3}/4540 = \qty{3104}{m}$।

\textbf{13.} $e^{-40/36.2} = 0.33124$, $e^{-360/36.2}/9 = \num{5.3e-6}$ और
$n = 5$ पद और भी छोटा है; कोष्ठक का मान है $1 - 0.81057\times0.33125 =
0.73150$, और $d = 2500 + 3104\times0.73150 = \qty{4771}{m}$।

\textbf{14.} $n = 3$ पद से, $3104\times0.81057\times\num{5.3e-6} =
\qty{0.013}{m}$: \qty{40}{Myr} पर $e^{-9t/\tau}$ का मान, $10^{-4}$ गुना,
$e^{-t/\tau}$ के, और केवल पहला मोड शेष है।

\textbf{15.} \qty{40}{Myr} पर, $4843 - 4771 = \qty{73}{m}$;
\qty{50}{Myr} पर $2500 + 370.5\sqrt{50} = 5120$ की तुलना में $2500 + 3104\times(1 -
0.81057\times0.25127) = 4972$: \qty{148}{m}। अंतर्वेशन से, $40 + 10\times
(100 - 73)/(148 - 73) = \qty{43.6}{Myr}$, लगभग \qty{44}{Myr} (पूरी श्रेणी
से 44.0 प्राप्त होता है)।

\textbf{16.} $q = (3.1\times1450/\num{95e3})[1 + 2(0.33124 + 0.01204 +
0.00005)] = 47.32\times1.68666 = \qty{79.8}{mW.m^{-2}}$, यह अर्ध-समष्टि के
समान है: \qty{40}{Myr} पर आधार का असर अभी
सतह तक नहीं पहुँचा है।

\textbf{17.} $2500 + 3104 = \qty{5604}{m}$ और $kT_m/L = \qty{47.3}{mW.m^{-2}}$।

\textbf{18.} $d(120) = 2500 + aL\bigl[1 - (8/\pi^2)e^{-120/\tau}\bigr]$, जहाँ
$a = \rho_m\alpha T_m/\bigl(2(\rho_m - \rho_w)\bigr) = \qty{32.673}{m/km}$ और
$\tau = L^2/(\pi^2\kappa)$।

\textbf{19.} $L = 3200/\bigl(32.673\,[1 - 0.81057e^{-120/\tau}]\bigr)$।
$L = \qty{100}{km}$ से आरंभ करके: $\tau = \qty{40.13}{Myr}$, कोष्ठक 0.95925, $L =
\qty{102.10}{km}$; फिर $\tau = 41.83$, $L = 102.67$; $\tau = 42.30$, $L =
102.82$; $\tau = 42.43$, $L = 102.87$; अंततः $L = \qty{102.9}{km}$ पर स्थिर हो जाता है।

\textbf{20.} $\tau = \qty{42.48}{Myr}$, $aL = \qty{3362}{m}$: $d(40) = 2500 +
3362\times(1 - 0.81057e^{-40/42.48}) = \qty{4799}{m}$, जिसकी तुलना में अर्ध-समष्टि के लिए
\qty{4843}{m} और मापी गई \qty{4800}{m} हैं: प्लेट दोनों आयुओं पर मेल खाती है।

\textbf{21.} $\tau = \qty{42.5}{Myr}$; अनंतस्थ गहराई $2500 + 3362 =
\qty{5862}{m}$।

\textbf{22.} $3362\times0.81057\,e^{-t/42.48} = 100$ के लिए $t = 42.48\ln27.25 =
\qty{140}{Myr}$।

\textbf{23.} काल्पनिक महासागर की पुरानी तली से मेल खाने वाली प्लेट
\textbf{\qty[text-series-to-math=true]{102.9}{km}} मोटी है।
\end{solution}
